La ley\par

\[
M_n
=
\left(\frac{a_n}{a_0}\right)^{-6}
\]

es quizá la más sencilla algebraicamente y una de las más profundas conceptualmente.\par

La dinámica de Perron produce dos comportamientos complementarios.\par

Por una parte, existe una expansión:\par

\[
V_n=\varphi^n.
\]

En tres dimensiones espaciales, la escala lineal asociada es la raíz cúbica:\par

\[
\frac{a_n}{a_0}
=
\varphi^{n/3}.
\]

Por otra parte, existe un modo estable de memoria:\par

\[
M_n=\varphi^{-2n}.
\]

La expansión crece con \(\varphi^n\). La memoria estable decrece con \(\varphi^{-2n}\). Al eliminar el nivel \(n\), aparece necesariamente\par

\[
M_n
=
\left(\frac{a_n}{a_0}\right)^{-6}.
\]

El exponente \(-6\) se deriva de dos hechos internos:\par

\begin{itemize}[leftmargin=*,itemsep=0.35em]
\item el espacio tiene dimensión \(3\), por lo que la escala lineal lleva exponente \(1/3\);
\item el modo estable es cuadrático, por lo que lleva exponente \(-2\).
\end{itemize}

La razón entre ambos es\par

\[
\frac{-2}{1/3}=-6.
\]

En Einstein–Cartan, la densidad de espín escala como un volumen inverso:\par

\[
s\sim a^{-3}.
\]

Su contribución cuadrática escala entonces como\par

\[
s^2\sim a^{-6}.
\]

La ley HMT reproduce exactamente ese exponente antes de introducir la realización Einstein–Cartan. Eso es lo verdaderamente sugestivo.\par

El modo de memoria de Perron se comporta como debe comportarse una densidad cuadrática de espín en una geometría con torsión.\par

La interpretación natural es que la torsión publica físicamente la memoria genealógica complementaria a la información métrica.\par

La curvatura describe cómo se deforma la geometría. La torsión describe cómo la geometría conserva una orientación de transporte adicional a las distancias. HMT posee desde el principio orientación, acarreo y memoria de ruta. Por eso la correspondencia con Einstein–Cartan es estructural: ambos marcos distribuyen la historia geométrica entre métrica y torsión.\par

Cuando el universo se expande, esa memoria se diluye como \(a^{-6}\). Se diluye más deprisa que una densidad de materia ordinaria o de radiación. Pero precisamente por eso puede ser dominante en la región de escala pequeña.\par

En un régimen primordial, una contribución torsional puede competir con la densidad ordinaria y modificar la singularidad. La posibilidad de un rebote depende del signo, de la amplitud, de la corriente de espín y de las condiciones de borde. HMT suministra primero la ley de escala exacta; la realización física fija después esas condiciones.\par

Eso ofrece una división del trabajo muy clara:\par

\begin{itemize}[leftmargin=*,itemsep=0.35em]
\item HMT genera la memoria estable y su exponente;
\item la realización Einstein–Cartan asigna esa memoria a una densidad cuadrática de espín;
\item la dinámica cosmológica fija amplitud, signo y condiciones físicas.
\end{itemize}

La relación con el principio de Mach aparece al juntar esta ley con la derivación de \(G\).\par

En la derivación gravitatoria, el acoplamiento local se selecciona mediante el cierre de un ciclo global. En la ley cosmológica, el estado local de torsión conserva la memoria de la expansión global. En ambos casos, la física local queda definida relacionalmente por la genealogía global.\par

La geometría local recuerda el estado global.\par

Éste es el sentido más concreto en que la construcción es machiana:\par

\begin{itemize}[leftmargin=*,itemsep=0.35em]
\item la inercia participa del reloj global;
\item la gravedad pertenece a la clausura de fase;
\item la torsión encarna la memoria de la historia;
\item la escala local codifica el nivel global de expansión.
\end{itemize}

El principio de Mach adquiere entonces una forma operatoria: las magnitudes locales son evaluaciones de una genealogía global, y los invariantes físicos son aquello que sobrevive al transporte entre ambas escalas.\par

La rama de periodicidad temporal de (108) fases y la rama de Perron expresan además dos aspectos complementarios de la cosmología.\par

La periodicidad temporal de (108) fases fija una escala de periodicidad y, mediante una realización de Sitter, una escala de curvatura. Es la cosmología del cierre: pregunta qué radio permite que el ciclo sea coherente.\par

Perron fija una tasa de expansión y una ley de dilución de memoria. Es la cosmología de la evolución: pregunta cómo cambia la genealogía al recorrer niveles.\par

Una rama describe el cierre de la geometría. La otra describe la persistencia de su memoria.\par

El calendario cronológico y la dinámica de Perron permanecen como realizaciones tipológicamente distintas y complementarias. Juntas producen una imagen especialmente rica: el universo posee una escala global de fase y una dinámica interna de memoria. Curvatura y torsión son dos publicaciones de esos dos aspectos.\par

\section{Confluencia de los lectores de vacío, acción, gravedad, área, mezcla y memoria}

Si se leen conjuntamente los nueve resultados, aparece una gramática común.\par

Cada lector publica selectivamente un invariante y transporta el resto como memoria complementaria.\par

El vacío conserva:\par

\begin{itemize}[leftmargin=*,itemsep=0.35em]
\item el producto como propagación;
\item el cociente como orientación.
\end{itemize}

La pareja \(G,\hbar\) conserva:\par

\begin{itemize}[leftmargin=*,itemsep=0.35em]
\item la acción como orientación;
\item el área como geometría común.
\end{itemize}

El cierre gravitatorio conserva:\par

\begin{itemize}[leftmargin=*,itemsep=0.35em]
\item la igualdad entre fase, Compton, gravedad y Planck.
\end{itemize}

El modo \(30\) conserva:\par

\begin{itemize}[leftmargin=*,itemsep=0.35em]
\item la cadencia común de vacío, criticidad y memoria.
\end{itemize}

El área y el Hamiltoniano modular conservan:\par

\begin{itemize}[leftmargin=*,itemsep=0.35em]
\item la cantidad total;
\item la procedencia sectorial.
\end{itemize}

El tránsito \(6\to8\) conserva:\par

\begin{itemize}[leftmargin=*,itemsep=0.35em]
\item el giro geométrico en Barbero;
\item su huella informacional en la entropía;
\item su orientación de mezcla en CKM.
\end{itemize}

El determinante angular conserva:\par

\begin{itemize}[leftmargin=*,itemsep=0.35em]
\item la parte par que puede convertirse en mezcla electrodébil.
\end{itemize}

La torsión de cuarto de giro conserva:\par

\begin{itemize}[leftmargin=*,itemsep=0.35em]
\item una estructura compleja que se publica como carácter aritmético.
\end{itemize}

La dinámica de Perron conserva:\par

\begin{itemize}[leftmargin=*,itemsep=0.35em]
\item la memoria estable de la expansión, que se realiza como torsión cosmológica.
\end{itemize}

Por eso las constantes se relacionan mediante una arquitectura común de transporte y conservan sus identidades funcionales. Son residuos diferentes de esa arquitectura.\par

La Holografía Modular Triádica aporta la genealogía.\par
La Mecánica Dimensional aporta la realización.\par
El número real aparece al final como valor terminal.\par

En esa visión, una constante fundamental es una frase completa acerca del universo. Dice:\par

\begin{itemize}[leftmargin=*,itemsep=0.35em]
\item qué cambió;
\item qué orientación se invirtió;
\item qué memoria cruzó el borde;
\item qué información permaneció interior y qué parte quedó publicada;
\item qué combinación permaneció invariante;
\item y en qué magnitud física se realizó finalmente esa invariancia.
\end{itemize}

Eso es lo que hace espectaculares estas derivaciones. Explican conjuntamente cuánto valen las constantes y por qué cada una desempeña exactamente su función estructural.\par
